\chapter{单片机用户界面} \label{ch:menu-system}

人机交互是嵌入式系统的重要环节。对于带显示屏和按键的单片机设备，菜单是最直观、最常用的交互形式。然而在单片机环境下实现菜单系统面临三个典型困难：其一，使用 \texttt{malloc} 与 \texttt{free} 动态分配内存会产生碎片，长期运行后可能耗尽内存；其二，C++ 的面向对象特性在部分编译环境下不够稳定，且很多单片机工程师只使用 C；其三，LVGL 等大型 GUI 库功能强大但代码量大、学习曲线陡峭，不适合资源有限的新手项目。第~\ref{ch:oop-in-c}~章介绍了用 C 语言实现面向对象特性的方法，本章的菜单系统正是这些方法的一个完整应用。

本章介绍的菜单系统采用“纯 C + 静态创建 + 缓冲区分配”的方案解决上述痛点，并通过函数指针在 C 语言中模拟“类与继承”，实现了跳转、退出、函数、数值修改、单选、多选等 6 类菜单项。项目同时在 PC 上使用 Win32 API（GDI）做了仿真，可以在电脑上先调试菜单逻辑，再移植到单片机。

\medskip \textbf{本章涉及到的实验内容：}
\begin{enumerate}
\item 菜单系统基础实验（菜单与菜单项的静态创建）；
\item 六类菜单项实验（跳转、退出、函数、数值、单选和多选）；
\item 数值修改模式实验（整数和浮点数、多级步进）；
\item 菜单系统 PC 仿真实验（Win32 API 窗口模拟 LCD）；
\item 单片机代码移植实验。
\end{enumerate}

% ---------------------------------------------------------------------
\section{设计思想}

% ---------------------------------------------------------------------
\subsection{文件架构}

菜单系统按职责分为三个文件（加上硬件移植层），如图~\ref{fig:menu-arch}~所示。

\begin{figure}[htbp]
\centering
    \begin{tikzpicture}[
        box/.style={draw=structurecolor, thick, rounded corners=2pt, align=center,
                    fill=structurecolor!8, minimum width=9.5cm, minimum height=0.85cm,
                    font=\small},
        arrow/.style={-stealth, thick}
    ]
        \node[box] (user) at (0,3.1) {menu\_user.c / menu\_user.h（用户创建菜单、主界面函数）};
        \node[box] (core) at (0,1.5) {menu.c / menu.h（菜单核心：创建、按键、显示、回调）};
        \node[box] (port) at (0,0)   {menu\_port.c / menu\_port.h（硬件接口：LCD、按键）};
        \node[box, minimum width=5.5cm] (hw)   at (-5,-1.5){硬件：LCD 液晶 / 矩阵键盘};

        \node[box, fill=second!8, draw=second] (sim) at (5,-1.5)
            {PC 仿真：win32\_sim.cpp（Win32 窗口模拟 LCD）、win\_main.c（主循环）};

        \draw[arrow] (user) -- (core);
        \draw[arrow] (core) -- (port);
        \draw[arrow] (port) -- (hw);
        \draw[arrow] (port) -- (sim);
    \end{tikzpicture}
\caption{菜单系统文件架构} \label{fig:menu-arch}
\end{figure}

各文件职责：

\begin{itemize}
\item \textbf{menu\_port}：定义菜单系统与硬件的接口（清屏、显示字符串、按键扫描等），移植时只需要修改这个文件；
\item \textbf{menu}：菜单系统的核心，负责菜单项的创建、菜单栈管理、按键分发与界面刷新；
\item \textbf{menu\_user}：用户代码，通过菜单 API 构建自己的菜单树，并编写主界面与功能界面；
\item \textbf{win32\_sim.cpp 与 win\_main.c}：PC 仿真层，用 Win32 API 窗口模拟 LCD，用键盘模拟按键，让菜单逻辑先在电脑上跑通。
\end{itemize}

% ---------------------------------------------------------------------
\subsection{用 C 模拟“类”}

菜单项千变万化（跳转项、数值项、单选、多选……），但它们的“行为”是统一的：绘制、进入、退出、修改。因此菜单系统借鉴面向对象思想，把每个菜单项看作一个“对象”，用函数指针表（相当于虚函数表）描述它的行为，用联合体（\texttt{union}）存放不同菜单项特有的数据。

\begin{lstlisting}[style=c]
// 菜单函数事件
typedef void (*UI_FuncEvent)(UI_MenuItem* sender, int16\_t inputKey);
typedef void (*UI_Event)(UI_MenuItem* sender);

struct _UI_MenuItem                  // 所有菜单项目的类
{
    const char* ItemTitle;           // 选项标题字符串
    UI_Menu* BaseMenu;               // 基于的父菜单

    UI_FuncEvent DrawDetail;         // 绘制该项的详情
    UI_FuncEvent Change;             // 修改事件
    UI_Event Enter;                  // 进入事件
    UI_Event Exit;                   // 退出事件

    union _ChildUnion
    {
        UI_Menu* Jump;               // 需要跳转的菜单
        UI_FuncEvent Func;           // 需要进入的函数
        uint32_t* BindPtr;           // 菜单变量绑定

        struct _ValueItem
        {
            uint32_t* BindPtr;       // 菜单变量绑定
            uint32_t Min;            // 修改下限
            uint32_t Max;            // 修改上限
            UI_MenuChangeValueType Type; // 变量修改类型
        } Value;
        struct _ChooseItem
        {
            uint32_t* BindPtr;       // 菜单变量绑定
            int32_t Establish;       // 选择成立值
        } Choose;
    } Union;
};
\end{lstlisting}

其中 \texttt{DrawDetail}、\texttt{Change}、\texttt{Enter} 与 \texttt{Exit} 四个函数指针就是“虚函数”：创建不同类型的菜单项时，只需要把对应的回调函数赋给这四个指针，调度器在运行时就只关心这四个接口，不需要关心菜单项的具体类型。

% ---------------------------------------------------------------------
\section{核心数据结构}

% ---------------------------------------------------------------------
\subsection{菜单与菜单项链表}

菜单（\texttt{UI\_Menu}）持有它的项目列表，项目列表用双向链表组织，便于上下移动；所有菜单又通过单向链表 \texttt{UI\_MenuListNode} 串成一个菜单列表。

\begin{lstlisting}[style=c]
struct _UI_MenuItemListNode           // 项目链表节点（双向）
{
    UI_MenuItem* Item;                // 当前项目
    UI_MenuItemListNode* pNext;       // 下一个项目
    UI_MenuItemListNode* pPrev;       // 上一个项目
    uint16_t Index;                   // 菜单所在的索引
};

struct _UI_Menu                       // 菜单
{
    const char* MenuTitle;            // 菜单标题
    UI_MenuItemListNode* MenuItemList;// 菜单项目列表
    UI_Driver* BaseUI;                // 基于的父界面

    UI_MenuItemListNode* StartItemNode;  // 列表开始显示的项目节点
    UI_MenuItemListNode* SelectItemNode; // 当前选中的项目节点

    uint16_t ItemCount;               // 项目个数
    uint16_t CurrentIndex;            // 当前项目索引
};

struct _UI_MenuListNode               // 菜单链表节点（单向）
{
    UI_Menu* Item;                    // 当前菜单
    UI_MenuListNode* pNext;           // 下一个菜单
};
\end{lstlisting}

\texttt{StartItemNode} 与 \texttt{SelectItemNode} 分别指向“当前屏幕第一行显示的项目”和“当前选中项目”，当项目总数超过屏幕显示行数时，选中光标移动会带动起始项目一起滚动，实现菜单翻页效果。

% ---------------------------------------------------------------------
\subsection{UI 驱动器}

\texttt{UI\_Driver} 是整个菜单系统的“总控制台”，包含菜单缓冲区、菜单栈和运行标志。

\begin{lstlisting}[style=c]
struct _UI_Driver                    // UI驱动器
{
    uint16_t MenuBufferTop;          // 菜单缓冲区已使用大小
    uint8_t MenuCount;               // 菜单数目
    uint8_t StackTop;                // 菜单栈顶
    uint8_t UpdateAllFlag;           // 全部刷新标志
    uint8_t UpdateValueFlag;         // 数值刷新标志
    uint8_t ValueFlag;               // 数值修改模式标志
    uint8_t FunctionRunFlag;         // 执行函数标志
    uint8_t LastFunctionRunFlag;     // 上次执行函数标志
    uint8_t ChangeUnit;              // 修改单位

    UI_FuncEvent MainUI;             // 主界面的函数
    UI_FuncEvent RunningFun;         // 正在运行的函数

    UI_MenuListNode* MenuList;       // 菜单列表（第0个为初始调用菜单）
    UI_Menu* MenuStack[UI_MENU_STACK_SIZE]; // 菜单栈
    uint8_t MenuBuffer[UI_MENU_BUFFER_SIZE]; // 菜单缓冲区
};
\end{lstlisting}

其中 \texttt{MenuStack} 是实现多级菜单的关键：进入子菜单时压栈，返回上一级时出栈； \texttt{MenuBuffer} 是 4KB 的静态缓冲区，所有菜单和菜单项都从这里分配，全程不依赖 \texttt{malloc}。

% ---------------------------------------------------------------------
\subsection{按键定义}

菜单系统定义了 6 个按键，如表~\ref{tab:menu-key}~所示（PC 仿真时直接映射到键盘）。

\begin{table}[htbp]
\centering \caption{菜单按键定义} \label{tab:menu-key}
    \begin{tabular}{cll}
        \toprule
        \textbf{宏} & \textbf{键值} & \textbf{功能} \\
        \midrule
        \texttt{UI\_KEY\_ENTER}    & Q & 进入或确认 \\
        \texttt{UI\_KEY\_EXIT}     & A & 退出或返回 \\
        \texttt{UI\_KEY\_ITEM\_UP} & W & 菜单项上移（或在修改模式切换单位） \\
        \texttt{UI\_KEY\_ITEM\_DOWN}& S & 菜单项下移（或在修改模式切换单位） \\
        \texttt{UI\_KEY\_VAL\_UP}   & E & 数值修改 +1 \\
        \texttt{UI\_KEY\_VAL\_DOWN} & D & 数值修改 -1 \\
        \bottomrule
    \end{tabular}
\end{table}

% ---------------------------------------------------------------------
\section{菜单的静态创建机制}

% ---------------------------------------------------------------------
\subsection{缓冲区分配器：UI\_BufferMalloc}

不使用 \texttt{malloc} 的原因在章首已经说明。菜单系统在 \texttt{UI\_Driver} 内预置了一块 4KB 静态缓冲区，所有菜单、菜单项、链表节点都从这块缓冲区顺序分配：

\begin{lstlisting}[style=c]
static void* UI_BufferMalloc(UI_Driver* pDirver, uint32_t size)
{
    // 判断分配空间是否溢出
    if (pDirver->MenuBufferTop + size > UI_MENU_BUFFER_SIZE) { return NULL; }
    // 分配新的空间并字节自增
    void* alloc = (void*)(pDirver->MenuBuffer + pDirver->MenuBufferTop);
    pDirver->MenuBufferTop += size;
    // 返回分配空间的首地址
    return alloc;
}
\end{lstlisting}

这种“顺序分配、永不释放”的分配器简单可靠：菜单结构在初始化阶段一次性创建完毕，之后不会再动态分配，因此不存在碎片问题；缓冲区大小由 \texttt{UI\_MENU\_BUFFER\_SIZE} 配置，初始化后可通过 \texttt{printf("缓存已使用\%d字节", MainUI.MenuBufferTop)} 观察实际使用量，从而确定缓冲区是否够用。

% ---------------------------------------------------------------------
\subsection{菜单与菜单项的添加}

\texttt{UI\_AddMenu} 负责创建菜单并挂到菜单列表尾部，\texttt{UI\_AddMenuItem} 负责把菜单项挂到指定菜单的项目链表尾部。以添加“跳转项”为例：

\begin{lstlisting}[style=c]
void UI_MenuAddJumpItem(UI_Menu* pMenu, const char* itemTitle, UI_Menu* jumpMenu)
{
    // 1. 分配空间
    UI_Driver* pDirver = pMenu->BaseUI;
    UI_MenuItem* newItem = (UI_MenuItem*)UI_BufferMalloc(pDirver, sizeof(UI_MenuItem));
    assert(newItem != NULL);
    // 2. 赋值相关信息
    newItem->ItemTitle = itemTitle;
    newItem->BaseMenu = pMenu;                   // 基于的菜单
    newItem->DrawDetail = UI_JumpItemDrawDetail; // 绘制函数
    newItem->Enter = UI_JumpItemEnter;           // 进入函数
    newItem->Exit = UI_BaseItemExit;             // 退出函数
    newItem->Change = UI_BaseItemChange;         // 修改函数
    newItem->Union.Jump = jumpMenu;              // 跳转目标菜单
    // 3. 添加到菜单
    UI_AddMenuItem(pMenu, (UI_MenuItem*)newItem);
}
\end{lstlisting}

可以看到，创建不同类型菜单项的代码结构完全一致：分配空间、绑定四个行为函数、填写联合体中的特有数据、挂入链表。这就是“工厂方法”的朴素实现。

% ---------------------------------------------------------------------
\subsection{六类菜单项}

系统共支持 6 类菜单项，如表~\ref{tab:menu-item}~所示。

\begin{table}[htbp]
\centering \caption{菜单项类型一览} \label{tab:menu-item}
    \begin{tabular}{lll}
        \toprule
        \textbf{类型} & \textbf{创建函数} & \textbf{行为} \\
        \midrule
        跳转项 & \texttt{UI\_MenuAddJumpItem}  & 进入后跳转到指定子菜单 \\
        退出项 & \texttt{UI\_MenuAddExitItem}  & 返回上一级菜单（或主界面） \\
        函数项 & \texttt{UI\_MenuAddFuncItem}  & 进入后运行指定函数界面 \\
        数值项 & \texttt{UI\_MenuAddValueItem\_Int} 与 \texttt{...\_Float} & 修改绑定的 int 与 float 变量 \\
        单选项目 & \texttt{UI\_MenuAddChooseItem} & 选中时把绑定变量置为指定值 \\
        多选项目 & \texttt{UI\_MenuAddSelectItem} & 在 ENABLE 与 DISABLE 之间切换 \\
        \bottomrule
    \end{tabular}
\end{table}

数值项是最常用的菜单项。进入数值项后系统进入“修改模式”，此时 \texttt{UI\_ValueItemChange} 按当前步进单位修改绑定变量，并限幅在 \texttt{Min~$\sim$~Max} 之间：

\begin{lstlisting}[style=c]
void UI_ValueItemChange(UI_MenuItem* pthis, int16_t delta)
{
    UI_Menu* pMenu = pthis->BaseMenu;
    UI_Driver* pDirver = pMenu->BaseUI;

    if (pDirver->ValueFlag == DISABLE) {
        // 若没有开启修改模式，先开启修改模式
        pDirver->ValueFlag = ENABLE;
        UI_ShowMode(pDirver);
        return;
    }

    // 根据不同的值类型进行修改
    switch (pthis->Union.Value.Type) {
    case UI_TYPE_INT: {
        int32_t* bindValPtr = UI_TypeTranPtr(pthis->Union.Value.BindPtr, int32_t);
        int32_t changeMin = UI_TypeTran(pthis->Union.Value.Min, int32_t);
        int32_t changeMax = UI_TypeTran(pthis->Union.Value.Max, int32_t);
        *bindValPtr = UI_LimitB(*bindValPtr + (int32_t)round(UI_GetUnit(pDirver) * delta),
                                changeMin, changeMax);
        break;
    }
    case UI_TYPE_FLOAT: {
        float* bindValPtr = UI_TypeTranPtr(pthis->Union.Value.BindPtr, float);
        float changeMin = UI_TypeTran(pthis->Union.Value.Min, float);
        float changeMax = UI_TypeTran(pthis->Union.Value.Max, float);
        *bindValPtr = UI_LimitA(*bindValPtr + UI_GetUnit(pDirver) * delta,
                                changeMin, changeMax);
        break;
    }
    }

    UI_ValueItemDrawDetail(pthis, pMenu->CurrentIndex + 1);
}
\end{lstlisting}

浮点数支持 4 档步进单位（见表~\ref{tab:menu-unit}），整数支持 2 档，在修改模式下再次按 “进入”键即可切换单位，便于快速调节到目标值。

\begin{table}[htbp]
\centering \caption{数值修改单位与步进值} \label{tab:menu-unit}
    \begin{tabular}{llll}
        \toprule
        \textbf{单位} & \textbf{显示} & \textbf{步进值} & \textbf{适用类型} \\
        \midrule
        \texttt{UI\_UNIT\_MA} & \texttt{<01.00F} & 1.00  & int 与 float \\
        \texttt{UI\_UNIT\_MB} & \texttt{<10.00F} & 10.00 & int 与 float \\
        \texttt{UI\_UNIT\_MC} & \texttt{<00.10F} & 0.10  & float \\
        \texttt{UI\_UNIT\_MD} & \texttt{<00.01F} & 0.01  & float \\
        \bottomrule
    \end{tabular}
\end{table}
% ---------------------------------------------------------------------
\section{菜单运行机制}

% ---------------------------------------------------------------------
\subsection{初始化与主循环}

使用菜单系统时，用户只需要在初始化阶段调用 \texttt{UI\_Init}，在主循环中不断调用 \texttt{UI\_While}：

\begin{lstlisting}[style=c]
int main(void)
{
    UI_Init();          // 硬件初始化 + 菜单创建

    while (1)
    {
        UI_FlushFlag = ENABLE;  // 置刷新标志位
        UI_While();             // 菜单运行
        Sleep(10);              // 控制帧率（单片机可改为短延时）
    }
}
\end{lstlisting}

\texttt{UI\_Init} 内部完成四步：硬件初始化、读取保存的变量、初始化驱动器（\texttt{UI\_StartInit}）、创建菜单树（\texttt{UI\_AddMenu} 等）、结束初始化（\texttt{UI\_EndInit}）。

\begin{lstlisting}[style=c]
void UI_Init(void)
{
    UI_HardwareInit();      // 1. LCD 等硬件初始化
    UI_Read();              // 2. 从EEPROM等存储读取变量
    UI_StartInit(&MainUI, UI_Main); // 3. 初始化驱动器

    // 4. 创建菜单树（见 3.6 节示例）
    UI_Menu* Menu_Root = UI_AddMenu(&MainUI, "主菜单");
    UI_Menu* Menu_Setting = UI_AddMenu(&MainUI, "数据设置");
    // ... 添加各类菜单项 ...

    UI_EndInit(&MainUI);    // 5. 结束初始化，进入主界面

    printf("缓存已使用%d字节！\n", MainUI.MenuBufferTop);
}
\end{lstlisting}

% ---------------------------------------------------------------------
\subsection{菜单栈与界面切换}

菜单系统使用“菜单栈”管理界面层级。进入子菜单（跳转项）时把新菜单压栈，返回上一级（退出项）时出栈，栈空时回到主界面函数。

\begin{lstlisting}[style=c]
// 跳转项：进入子菜单
void UI_JumpItemEnter(UI_MenuItem* pthis)
{
    UI_Driver* pDirver = pthis->BaseMenu->BaseUI;
    UI_Menu* newMenu = pthis->Union.Jump;
    // 设置初始选项
    newMenu->StartItemNode = newMenu->MenuItemList;
    newMenu->SelectItemNode = newMenu->MenuItemList;
    // 赋值到菜单栈
    pDirver->MenuStack[pDirver->StackTop++] = newMenu;
    pDirver->UpdateAllFlag = ENABLE; // 要求整屏刷新
}

// 退出项：返回上一级
void UI_BaseItemExit(UI_MenuItem* pthis)
{
    UI_Driver* pDirver = pthis->BaseMenu->BaseUI;
    if (pDirver->StackTop == 0) return;   // 栈空则忽略

    pDirver->StackTop--;                  // 出栈
    pDirver->UpdateAllFlag = ENABLE;

    if (pDirver->StackTop == 0)
    {
        UI_Clear();
        pDirver->RunningFun = pDirver->MainUI; // 回到主界面
        pDirver->FunctionRunFlag = ENABLE;
    }
}
\end{lstlisting}

% ---------------------------------------------------------------------
\subsection{按键处理：UI\_KeyPress}

\texttt{UI\_KeyPress} 是菜单运行的“大脑”。它先通过 \texttt{UI\_KeyScan} 检查是否有按键，然后按“进入、退出、数值修改、项目移动”的优先级分发按键事件：

\begin{lstlisting}[style=c]
static void UI_KeyPress(UI_Driver* pDirver)
{
    UI_Menu* currentMenu = pDirver->MenuStack[pDirver->StackTop - 1];
    UI_MenuItem* selectItem = currentMenu->SelectItemNode->Item;

    if (UI_KeyScan())   // 有按键
    {
        int8_t item = UI_GetKeyEvent(UI_MENU_EVENT_ITEM);
        int8_t value = UI_GetKeyEvent(UI_MENU_EVENT_VALUE);

        // 1. 进入
        if (UI_GetKeyEvent(UI_MENU_EVENT_ENTER) == ENABLE) {
            selectItem->Enter(selectItem);
        }
        // 2. 退出
        else if (UI_GetKeyEvent(UI_MENU_EVENT_EXIT) == ENABLE) {
            selectItem->Exit(selectItem);
        }
        // 3. 数值修改
        else if (value != 0) {
            if (UI_IsChangeMode(pDirver))
                selectItem->Change(selectItem, value);
        }
        // 4. 项目移动
        else {
            // ... 在项目链表上前后移动 SelectItemNode，
            //     超过屏幕范围时带动 StartItemNode 一起滚动 ...
            UI_ShowChooseList(currentMenu);   // 刷新选中光标
        }
    }
}
\end{lstlisting}

按键事件本身由移植层通过 \texttt{UI\_GetKeyEvent} 归一化为四类：进入、退出、项目偏移、数值偏移，菜单核心不需要关心具体按键的硬件细节。

% ---------------------------------------------------------------------
\subsection{显示刷新：UI\_ShowItems}

\texttt{UI\_ShowItems} 负责把当前菜单绘制到屏幕上：第一行显示菜单标题（含模式与单位），其余行从 \texttt{StartItemNode} 开始逐个调用菜单项的 \texttt{DrawDetail} 绘制。

\begin{lstlisting}[style=c]
static void UI_ShowItems(UI_Driver* pDirver)
{
    if (pDirver->StackTop == 0) { return; }
    UI_Menu* currentMenu = pDirver->MenuStack[pDirver->StackTop - 1];
    UI_MenuItem* selectItem = currentMenu->SelectItemNode->Item;

    if (pDirver->UpdateAllFlag == ENABLE || pDirver->UpdateValueFlag == ENABLE)
    {
        UI_MenuItemListNode* currentItemNode = currentMenu->StartItemNode;

        if (pDirver->UpdateAllFlag == ENABLE)
        {
            UI_Clear();
            UI_ShowChooseList(currentMenu);   // 绘制选中光标
            UI_ShowTitle(pDirver, currentMenu); // 绘制标题行
        }

        for (uint16_t line = 1; line < UI_MAX_CHAR_LINE && currentItemNode != NULL; line++)
        {
            currentItemNode->Item->DrawDetail(currentItemNode->Item, line);
            if (pDirver->UpdateAllFlag == ENABLE) { UI_DrawTitle(currentItemNode->Item, line); }
            currentItemNode = currentItemNode->pNext;
        }
        pDirver->UpdateAllFlag = DISABLE;
        pDirver->UpdateValueFlag = DISABLE;
    }
}
\end{lstlisting}

只有当 \texttt{UpdateAllFlag}（整屏刷新）或 \texttt{UpdateValueFlag}（数值刷新）置位时才真正重绘，平时直接跳过，从而避免频繁操作 LCD 导致闪烁。

% ---------------------------------------------------------------------
\section{移植接口}

把菜单系统移植到单片机上，只需要实现 \texttt{menu\_port.c} 中的接口，见表~\ref{tab:menu-port}~。

\begin{table}[htbp]
\centering \caption{菜单系统移植接口} \label{tab:menu-port}
    \begin{tabular}{lll}
        \toprule
        \textbf{接口} & \textbf{功能} & \textbf{说明} \\
        \midrule
        \texttt{UI\_HardwareInit()} & 硬件初始化 & 初始化 LCD、按键等 \\
        \texttt{UI\_Clear()}        & 清屏 & 调用 LCD 清屏函数 \\
        \texttt{UI\_ShowString(x,y,str)} & 显示字符串 & 调 LCD 字符串显示 \\
        \texttt{UI\_KeyScan()}      & 按键扫描 & 有键按下返回 ENABLE \\
        \texttt{UI\_GetKeyValue()}  & 获取键值 & 返回当前按键码 \\
        \texttt{UI\_GetKeyEvent(ev)}& 归一化事件 & 映射为进入、退出、移动和修改 \\
        \texttt{UI\_Read()} 与 \texttt{UI\_Save(val)} & 变量读写 & EEPROM 保存菜单参数 \\
        \bottomrule
    \end{tabular}
\end{table}

以按键接口为例，移植层只需把物理按键归一化为菜单系统的事件：

\begin{lstlisting}[style=c]
int8_t UI_GetKeyEvent(MenuEvent menuEvent)
{
    switch (menuEvent) {
    case UI_MENU_EVENT_ENTER:   // 进入事件
        if (Key_GetNum() == UI_KEY_ENTER) { Key_Clear(); return ENABLE; }
        return DISABLE;
    case UI_MENU_EVENT_EXIT:    // 退出事件
        if (Key_GetNum() == UI_KEY_EXIT) { Key_Clear(); return ENABLE; }
        return DISABLE;
    case UI_MENU_EVENT_ITEM:    // 项目移动
        if (Key_GetNum() == UI_KEY_ITEM_DOWN) { Key_Clear(); return 1; }
        if (Key_GetNum() == UI_KEY_ITEM_UP)   { Key_Clear(); return -1; }
        return 0;
    case UI_MENU_EVENT_VALUE:   // 数值修改
        if (Key_GetNum() == UI_KEY_VAL_UP)   { Key_Clear(); return 1; }
        if (Key_GetNum() == UI_KEY_VAL_DOWN) { Key_Clear(); return -1; }
        return 0;
    }
    return DISABLE;
}
\end{lstlisting}

% ---------------------------------------------------------------------
\section{用户菜单创建示例}

以项目自带的 \texttt{menu\_user.c} 为例，一个包含“数据设置、单选、多选、函数项目”的完整菜单树如图~\ref{fig:menu-tree}~所示。

\begin{figure}[htbp]
\centering
    \begin{tikzpicture}[
        box/.style={draw=structurecolor, thick, rounded corners=2pt, align=center,
                    fill=structurecolor!8, font=\small, inner sep=3pt},
        arrow/.style={-stealth, thick}
    ]
        \node[box] (root) at (0,0) {主菜单};
        \node[box] (set)  at (-5.2,-1.8) {数据设置};
        \node[box] (radio) at (-1.8,-1.8) {单选按钮};
        \node[box] (check) at (1.8,-1.8) {多选按钮};
        \node[box] (func) at (5.2,-1.8) {函数项目};

        \node[box, fill=second!8, draw=second] (set2) at (-5.2,-3.8) {整数变量 A\\浮点数 B};
        \node[box, fill=second!8, draw=second] (radio2) at (-1.8,-3.8) {选项0~9};
        \node[box, fill=second!8, draw=second] (check2) at (1.8,-3.8) {选项0~9};
        \node[box, fill=second!8, draw=second] (func2) at (5.2,-3.8) {智能车图像\\陀螺仪解算};

        \draw[arrow] (root) -- (set);
        \draw[arrow] (root) -- (radio);
        \draw[arrow] (root) -- (check);
        \draw[arrow] (root) -- (func);
        \draw[arrow] (set) -- (set2);
        \draw[arrow] (radio) -- (radio2);
        \draw[arrow] (check) -- (check2);
        \draw[arrow] (func) -- (func2);
    \end{tikzpicture}
\caption{示例菜单结构} \label{fig:menu-tree}
\end{figure}

创建这段菜单树的代码非常直观，每类菜单项一行即可：

\begin{lstlisting}[style=c]
int32_t A = 0;      // 绑定的整数变量
float  B = 0;       // 绑定的浮点变量
int32_t Radio = 0;  // 单选绑定变量
int32_t Check[10] = { 0 }; // 多选绑定数组

void UI_Init(void)
{
    UI_HardwareInit();
    UI_Read();
    UI_StartInit(&MainUI, UI_Main);

    /* 创建菜单 */
    UI_Menu* Menu_Root     = UI_AddMenu(&MainUI, "主菜单");
    UI_Menu* Menu_Setting  = UI_AddMenu(&MainUI, "数据设置");
    UI_Menu* Menu_RadioBox = UI_AddMenu(&MainUI, "单选按钮");
    UI_Menu* Menu_CheckBox = UI_AddMenu(&MainUI, "多选按钮");
    UI_Menu* Menu_Function = UI_AddMenu(&MainUI, "函数项目");

    /* 主菜单：退出 + 4个跳转项 */
    UI_MenuAddExitItem(Menu_Root, "返回上一级");
    UI_MenuAddJumpItem(Menu_Root, "数据设置", Menu_Setting);
    UI_MenuAddJumpItem(Menu_Root, "单选按钮", Menu_RadioBox);
    UI_MenuAddJumpItem(Menu_Root, "多选按钮", Menu_CheckBox);
    UI_MenuAddJumpItem(Menu_Root, "函数项目", Menu_Function);

    /* 数据设置：整数变量 + 浮点数变量 */
    UI_MenuAddExitItem(Menu_Setting, "返回上一级");
    UI_MenuAddValueItem_Int(Menu_Setting, "整数变量", &A, 0, 100);
    UI_MenuAddValueItem_Float(Menu_Setting, "浮点数变量", &B, 0, 100);

    /* 单选按钮：10个选项 */
    UI_MenuAddExitItem(Menu_RadioBox, "返回上一级");
    for (int i = 0; i < 10; i++) {
        UI_MenuAddChooseItem(Menu_RadioBox, "选项%d", &Radio, i);
    }

    /* 多选按钮：10个开关 */
    UI_MenuAddExitItem(Menu_CheckBox, "返回上一级");
    for (int i = 0; i < 10; i++) {
        UI_MenuAddSelectItem(Menu_CheckBox, "选项%d", &Check[i]);
    }

    /* 函数项目 */
    UI_MenuAddExitItem(Menu_Function, "返回上一级");
    UI_MenuAddFuncItem(Menu_Function, "智能车图像", UI_Image);
    UI_MenuAddFuncItem(Menu_Function, "陀螺仪解算", UI_MPU);

    UI_EndInit(&MainUI);
}
\end{lstlisting}

主界面函数 \texttt{UI\_Main} 与功能界面函数（\texttt{UI\_Image}、\texttt{UI\_MPU}）使用相同的签名：\texttt{void func(UI\_MenuItem* sender, int16\_t inputKey)}。系统进入功能界面后会把按键逐个传给该函数，函数内部按 \texttt{inputKey} 决定行为；按下退出键时调用 \texttt{UI\_StopRunFunc} 结束功能界面，返回菜单模式。

% ---------------------------------------------------------------------
\section{页面跳转框架}

菜单系统里的“菜单栈”能管理父子界面，但栈会带来“返回路径不可变”的限制：一个页面只能沿进入顺序原路返回。对于主页、菜单页、曲线页这类平级页面，更简单的做法是\textbf{不用栈，用显式跳转}：谁想进入下一个页面，就由谁直接调用跳转函数。

\begin{figure}[htbp]
	\centering
	\begin{tikzpicture}[
		state/.style={rectangle, draw=structurecolor, thick, rounded corners=5pt,
			minimum width=2.5cm, minimum height=1cm, align=center, font=\small,
			fill=structurecolor!8},
		arrow/.style={-stealth, thick, structurecolor},
		label/.style={font=\footnotesize, auto, sloped}
		]
		
		% 三个页面状态
		\node[state] (home) at (0,0) {主页};
		\node[state] (menu) at (4,0) {菜单页};
		\node[state] (curve) at (8,0) {曲线页};
		
		% 跳转关系（带按键标注）
		\draw[arrow] (home) to[out=0,in=90] node[label, above] {ENTER} (menu);
		\draw[arrow] (menu) to[out=0,in=90] node[label, above] {DOWN} (curve);
		\draw[arrow] (curve) to[out=180,in=-90] node[label, below] {MODE} (menu);
		\draw[arrow] (menu) to[out=180,in=-90] node[label, below] {MODE} (home);
		
	\end{tikzpicture}
	\caption{页面跳转框架状态转换图}
	\label{fig:page-state}
\end{figure}

\subsection{页面导航数据结构}

\begin{lstlisting}[style=c]
typedef void (*UI_PageEntry_T)(void *param);

typedef struct {
    UI_PageEntry_T currentPage;  /* 当前页面入口函数 */
    void *currentParam;          /* 当前页面参数 */
} UI_Navigator_T;

static UI_Navigator_T s_Navigator = { 0 };
\end{lstlisting}

\texttt{currentPage} 保存当前页面函数指针，\texttt{currentParam} 保存进入页面时传入的参数。页面切换本质上只是“把这两个成员替换成新页面，并调用新页面入口”。

\subsection{页面跳转 API}

\begin{lstlisting}[style=c]
void GUI_JumpPage(UI_PageEntry_T page, void *param)
{
    s_Navigator.currentPage = page;
    s_Navigator.currentParam = param;

    page(param);   /* 立即进入新页面 */
}

void GUI_ReturnPage(void *param)
{
    /* 返回由调用方显式指定，例如回到主页面 */
    GUI_JumpPage(GUI_HomePage, param);
}
\end{lstlisting}

如果需要在跳转前执行旧页面的退出动作，可在 \texttt{GUI\_JumpPage} 中先调用 \texttt{currentPage->Exit} 或保存退出回调；最小框架中省略退出步骤。

\subsection{使用示例：主页 → 菜单页 → 曲线页}

下面三个页面均为普通函数，页面之间的跳转关系由代码显式写出：

\begin{lstlisting}[style=c]
void GUI_HomePage(void *param)
{
    /* 绘制主页 */
    if (Key_IsAction(KEY_ENTER)) {
        GUI_JumpPage(GUI_MenuPage, NULL);   /* 主页进入菜单页 */
    }
}

void GUI_MenuPage(void *param)
{
    /* 绘制菜单页 */
    if (Key_IsAction(KEY_DOWN)) {
        GUI_JumpPage(GUI_CurvePage, NULL);  /* 菜单页进入曲线页 */
    }
    if (Key_IsAction(KEY_MODE)) {
        GUI_JumpPage(GUI_HomePage, NULL);   /* 显式返回主页 */
    }
}

void GUI_CurvePage(void *param)
{
    /* 绘制曲线页 */
    if (Key_IsAction(KEY_MODE)) {
        GUI_JumpPage(GUI_MenuPage, NULL);   /* 显式返回菜单页 */
    }
}
\end{lstlisting}

每个页面都不依赖“上一个页面是谁”，因此可以自由组成主页→菜单页→曲线页、曲线页→主页等任意跳转路径。

\subsection{栈式导航与显式跳转对比}

\begin{table}[htbp]
\centering \caption{菜单栈式导航与显式页面跳转的对比} \label{tab:nav-stack-vs-explicit}
    \begin{tabular}{lll}
        \toprule
        \textbf{对比项} & \textbf{菜单栈导航} & \textbf{显式页面跳转} \\
        \midrule
        返回方式 & 只能按入栈顺序返回 & 可在任意页面指定目标 \\
        页面关系 & 父子层级固定 & 平级页面自由跳转 \\
        数据结构 & 需要页面栈与深度限制 & 只需要当前页面指针 \\
        适合场景 & 层级菜单、设置向导 & 主页与功能页切换 \\
        缺点 & 无法“跨级跳转” & 返回路径需要每页自己写 \\
        \bottomrule
    \end{tabular}
\end{table}

一句话总结：\textbf{页面是状态，跳转是有向边。}把页面看成节点，显式跳转就是在状态图上沿着“有向边”移动；相比维护一个“只能往回走”的栈，显式跳转更适合菜单与自定义页面、曲线页面之间的自由切换。

% ---------------------------------------------------------------------
\section{PC 仿真实现}

为了让菜单逻辑不依赖真实硬件即可调试，项目用 Win32 API（GDI）在 PC 上做了仿真：

\begin{itemize}
\item \textbf{win32\_sim.cpp}：用 \texttt{CreateWindow} 创建一个 240$\times$300 的窗口模拟 LCD，把 \texttt{LCD\_Init}、\texttt{LCD\_Clear} 与 \texttt{LCD\_ShowString} 映射为 Win32 GDI 绘图函数；
\item \textbf{按键映射}：用 \texttt{GetKeyState} 检测键盘上的 Q、A、W、S、E 和 D，映射为 6 个菜单按键；
\item \textbf{主循环}：创建窗口并进入消息循环，以 10ms 为周期调用 \texttt{UI\_While()}。
\end{itemize}

页面跳转框架的\textbf{独立 PC 仿真}位于 \texttt{Chapter-ui-system/页面跳转PC仿真/}，它在独立窗口中演示主页、菜单页和曲线页跳转，并额外提供 \texttt{LCD\_DrawPoint} 与 \texttt{LCD\_ShowRGB565Image} 绘制接口。

\begin{lstlisting}[style=c]
// win32_sim.cpp：LCD 接口的 Win32 API 实现（对接 hardware.h）
static HWND s_hWnd = NULL;
static HDC  s_hdc = NULL;

extern "C" void LCD_Init(void)
{
    /* 由本文件的 WinMain 创建窗口，这里只保存窗口 DC 与字体 */
    s_hdc = GetDC(s_hWnd);
    SetBkMode(s_hdc, TRANSPARENT);
    SetTextColor(s_hdc, RGB(0, 0, 0));
    SelectObject(s_hdc, CreateFontA(UI_CHAR_SIZE_HEIGHT, 0,
                 0, 0, FW_NORMAL, 0, 0, 0,
                 DEFAULT_CHARSET, 0, 0, CLEARTYPE_QUALITY, 0, "宋体"));
}

extern "C" void LCD_Clear(void)
{
    RECT rect = { 0, 0, UI_SIZE_WIDTH, UI_SIZE_HEIGHT };
    FillRect(s_hdc, &rect, (HBRUSH)GetStockObject(WHITE_BRUSH));
}

extern "C" void LCD_ShowString(uint16_t x, uint16_t y, const char *str)
{
    int length = (int)strlen(str);
    RECT rect = {
        x, y,
        x + length * UI_CHAR_SIZE_WIDTH,
        y + UI_CHAR_SIZE_HEIGHT
    };

    /* 先填充字符区域，避免旧字符残留造成文字重叠 */
    FillRect(s_hdc, &rect, (HBRUSH)GetStockObject(WHITE_BRUSH));
    TextOutA(s_hdc, x, y, str, length);
}
\end{lstlisting}

仿真时菜单核心（\texttt{menu.c} 与 \texttt{menu.h}）与真实硬件上的代码完全一致，只有移植层（\texttt{menu\_port}）被替换为 Win32 API 实现，相当于在 PC 上用 WinAPI“对接”了 HAL 层。因此 PC 上验证通过的菜单逻辑可以直接搬到单片机，真正做到“一次编写，两端复用”。

\begin{anybox}{PC 仿真一致性}
PC 仿真里的 \texttt{LCD\_ShowString} 也应先填充字符区域再绘制文字，否则滚动或局部刷新时会出现与真实 LCD 相同的字符残影。
\end{anybox}


\subsection{与 LVGL 现成框架的对比}

本章介绍的菜单系统与 LVGL（Light and Versatile Graphics Library）等成熟的嵌入式 GUI 框架，虽然在“显示屏 + 按键”人机交互的宏观目标上一致，但在设计哲学、资源占用和适用场景上有根本性差异。把两者放在一起对比，有助于读者理解“什么时候应该自己写菜单，什么时候应该用 LVGL”。

LVGL 是目前嵌入式领域最流行的开源 GUI 库之一，支持按钮、滑块、列表、图表、键盘、动画等丰富的控件，并提供了跨平台的显示驱动与输入设备抽象层。其核心定位是“为嵌入式系统提供接近智能手机体验的图形界面”。表~\ref{tab:menu-vs-lvgl}~从多个维度进行了对比。

\begin{table}[htbp]
	\centering
	\small
	\caption{本菜单系统与 LVGL 的对比}
	\label{tab:menu-vs-lvgl}
	\begin{tabularx}{\textwidth}{l X X}
		\toprule
		\textbf{对比项} & \textbf{本章菜单系统} & \textbf{LVGL（v8/v9）} \\
		\midrule
		设计定位 & 轻量级按键式菜单框架 & 完整的嵌入式 GUI 图形库 \\
		UI 控件类型 & 6 类菜单项（跳转/退出/函数/数值/单选/多选） & 30+ 控件（按钮、滑块、图表、开关等） \\
		图形能力 & 字符 + 光标（纯文本界面） & 位图、抗锯齿、透明度、圆角、阴影、动画 \\
		输入方式 & 5~6 个离散按键 & 触摸屏、鼠标、键盘、编码器、按键（全支持） \\
		界面布局 & 固定行列表单 & 灵活布局（flex、grid、绝对定位） \\
		内存占用 & 数 KB（RAM：TCB + 菜单树 + 栈） & 数十到数百 KB（帧缓冲 + 控件对象 + 样式 + 动画） \\
		Flash 占用 & 约 10--20 KB（仅菜单核心逻辑） & 约 100--500 KB（核心库 + 控件 + 字体 + 驱动） \\
		运行机制 & 主循环轮询按键，协作式刷新 & 事件驱动 + 定时器刷新 \\
		用户代码风格 & 手动创建菜单树，函数指针绑定行为 & 创建控件对象，注册事件回调 \\
		学习曲线 & 低（看懂结构体即可上手） & 中高（需理解对象树、样式继承、事件系统等） \\
		调试难度 & 低（单步跟踪菜单栈与状态即可） & 中（需配合 SquareLine Studio） \\
		移植难度 & 低（只需实现 LCD 绘字符串 + 按键扫描） & 中（需实现显示驱动和输入驱动） \\
		适用 MCU & Cortex-M0/M3，RAM $>$ 4KB 即可 & 建议 Cortex-M4/M7，RAM $>$ 32KB，Flash $>$ 256KB \\
		开发工具 & 无，纯手写代码 & SquareLine Studio、EEZ Studio（UI 拖拽生成代码） \\
		开源协议 & 示例代码（可自由修改使用） & MIT（自由使用，需保留版权声明） \\
		\bottomrule
	\end{tabularx}
\end{table}

从表中可以提炼出三条核心差异：

\begin{enumerate}
	\item \textbf{抽象层次不同}：LVGL 提供的是“控件”级别的抽象——开发者创建按钮、滑块、列表，并设置它们的样式和事件回调；本章菜单系统的抽象层次是“菜单项”——每一行要么是跳转、要么是数值修改、要么是函数执行。前者适用于图形化触摸界面，后者适用于物理按键 + 字符屏的文本菜单。
	
	\item \textbf{资源开销的鸿沟}：一个最小的 LVGL 工程（仅核心库 + 帧缓冲）通常需要 32KB RAM 和 128KB Flash，驱动 320x240 分辨率的屏幕时帧缓冲本身就要 150KB 以上。如果再加上字体、控件和样式，RAM 占用很容易突破 256KB。而本章菜单系统在相同屏幕下只需数 KB RAM 和约 10--20KB Flash，可以运行在 STM32F103C8T6（20KB RAM/64KB Flash）等入门级 MCU 上。
	
	\item \textbf{开发范式的取舍}：LVGL 的控件 + 样式 + 事件模型提供了极高的 UI 表现力，但开发者需要理解对象树的生命周期管理、样式继承链、动画和定时器机制，入门门槛较高；本章菜单系统的“创建菜单树 + 绑定函数指针”范式更接近嵌入式工程师熟悉的“表驱动”风格，代码平铺直叙，不依赖框架的运行时规则。
\end{enumerate}

那么如何选择？以下原则可作为参考：

\begin{itemize}
	\item \textbf{使用 LVGL 的典型场景}：
	\begin{itemize}
		\item 硬件配备触摸屏，或需要滑动、拖动等手势交互；
		\item UI 需要包含图表、滑块、下拉列表等复杂控件；
		\item 产品或项目要求“现代化”的视觉风格和流畅动画；
		\item MCU 资源充裕（Cortex-M4/M7 以上，RAM $\geq$ 64KB，Flash $\geq$ 512KB）；
		\item 团队有精力投入图形界面设计，并愿意使用 SquareLine Studio 等辅助工具。
	\end{itemize}
	
	\item \textbf{使用本章菜单系统的典型场景}：
	\begin{itemize}
		\item 硬件只有物理按键（如编码器、五向键、独立按键），无触摸屏；
		\item 屏幕以字符显示为主（如 128x64 OLED、1602 LCD），或仅有少量图形绘制；
		\item UI 逻辑本质上是“列表选择 + 参数修改”，不需要复杂的布局变化；
		\item MCU 资源紧张（RAM $<$ 16KB），或 Flash 已分配给其他核心功能；
		\item 希望 UI 逻辑清晰可单步调试，不想引入大型第三方库增加不确定性。
	\end{itemize}
\end{itemize}

最后需要强调一点：\textbf{LVGL 和本菜单系统不是二选一的互斥关系}。一个实际产品中，完全可以在 LVGL 的某个界面中嵌入一个“参数列表”风格的菜单页，或在菜单系统的某个函数项中调用 LVGL 绘制的曲线画面。两者既可以共存，也可以分层：LVGL 负责展示层（仪表盘、曲线、状态图标），本菜单系统负责参数配置与设备控制的逻辑流程。读者应把二者视为工具链上的不同选项，而非“谁替代谁”的竞争关系。

% ---------------------------------------------------------------------
\section{单片机代码移植实现}

\begin{figure}[htbp]
\centering
    % TODO: 在此插入第6章菜单系统移植实测截图
\caption{第6章实验记录图（待补充）} \label{fig:ch6-placeholder}
\end{figure}

% ---------------------------------------------------------------------
\begin{problemset}
\item 菜单项用联合体复用“跳转/函数/数值/单选/多选”数据时，靠哪个字段判断当前应读取哪一成员？
\item 在静态内存池上创建菜单树有什么好处？如果不小心超出 \texttt{UI\_MENU\_BUFFER\_SIZE}，应该如何处理？
\item 页面栈导航与显式页面跳转分别适合什么场景？画出主页→菜单页→曲线页两种导航方式的区别。
\item 绘制文字前为什么要先“填充一块背景区域”？直接逐字符叠加显示会发生什么问题？
\item PC 仿真器把 \texttt{LCD\_DrawPoint}、\texttt{LCD\_ShowRGB565Image} 映射为 Win32 API，这种移植层隔离的价值是什么？
\end{problemset}